\chapter{Kidney Stone Treatment}
\index{Kidney Stone Treatment}

\section*{Definition and Overview}
Kidney stone treatment depends on the size, position, and type of the stone, and on whether it is causing pain, obstruction, or infection. Many small stones pass on their own with fluids, pain relief, and sometimes medicines that relax the ureter. Stones that are too large to pass, cause persistent pain, block urine flow, or are infected need active treatment, including shock wave lithotripsy, ureteroscopy, or surgery. In addition to treating the stone itself, treatment includes preventing future stones with fluids, dietary change, and, in some cases, medicines. An obstructed or infected kidney stone is an emergency.

\section*{Epidemiology}
Around one in ten people develops a kidney stone, and treatment is needed for a substantial proportion. Most stones under 5 millimetres pass spontaneously, while most over 10 millimetres require intervention. Ureteroscopy has become the most common surgical treatment, replacing open surgery, and lithotripsy is widely available. Emergency presentations with obstructed or infected stones are common and can be life-threatening. Recurrence affects around half of people within a decade, making prevention a key part of treatment.

\section*{Aetiology and Pathophysiology}
Stones form when urine is concentrated with stone-forming substances, and their behaviour depends on size and site. Small stones in the kidney may be silent, while stones entering the ureter cause colicky pain and can block the flow of urine. Obstruction raises pressure in the kidney, damaging it over time, and an obstructed kidney with infection can develop into a urological emergency with sepsis. The choice of treatment matches the stone: observation for small, uncomplicated stones; lithotripsy for stones in the kidney or upper ureter; ureteroscopy for stones in the ureter; and surgery for very large stones. Preventing recurrence requires understanding the stone type and the underlying metabolic factors.

\section*{Clinical Features}
\begin{itemize}
    \item Small stones: often no treatment beyond fluids and pain relief, with the stone passing in days to weeks
    \item Pain relief: strong analgesics for colic, sometimes with medicines that relax the ureter
    \item Lithotripsy: shock waves from outside the body break stones into passable fragments
    \item Ureteroscopy: a telescope passed through the bladder into the ureter to remove or fragment the stone
    \item Surgery: keyhole or open removal for very large stones
    \item Stenting: a temporary tube to drain an obstructed kidney
    \item Emergency treatment: urgent drainage and antibiotics for infected obstruction
    \item Prevention: fluids, diet, and medicines to prevent recurrence
\end{itemize}

\section*{Differential Diagnosis}
The choice of treatment is individualised.
\begin{itemize}
    \item \textbf{Observation:} for small, painless stones
    \item \textbf{Medical therapy:} fluids, analgesics, and alpha-blockers to help passage
    \item \textbf{Shock wave lithotripsy:} for kidney and upper ureter stones
    \item \textbf{Ureteroscopy:} for ureteric stones and some kidney stones
    \item \textbf{Percutaneous nephrolithotomy:} keyhole removal of very large stones
    \item \textbf{Emergency drainage:} for obstruction, infection, or sepsis
    \item \textbf{Metabolic treatment:} prevention for recurrent stones
\end{itemize}

\section*{When to Seek Medical Care}
People with kidney stones should be managed with urology guidance, and any stone with severe pain, fever, or inability to pass urine needs urgent care. Stones that do not pass within a few weeks, or that grow, warrant active treatment. Anyone with recurrent stones should have metabolic assessment.

\textbf{Contact your local emergency services for:}
\begin{itemize}
    \item Severe flank pain with fever, shaking, or confusion, which may indicate sepsis
    \item Inability to pass urine
    \item Vomiting with an inability to keep fluids down
    \item Severe pain not relieved by prescribed medicines
    \item Collapse, dizziness, or severe illness
\end{itemize}

\section*{Diagnosis and Investigations}
\begin{itemize}
    \item \textbf{Imaging:} CT, ultrasound, or X-ray to locate and size the stone
    \item \textbf{Urine tests:} infection, blood, and crystals
    \item \textbf{Blood tests:} kidney function, calcium, and inflammatory markers
    \item \textbf{Stone analysis:} of stones passed or removed
    \item \textbf{24-hour urine collection:} for recurrent stones
    \item \textbf{Follow-up imaging:} confirming passage or resolution
    \item \textbf{Urology review:} planning the best treatment
\end{itemize}

\section*{Management}
\begin{itemize}
    \item \textbf{Conservative care:} fluids, analgesia, and time for small stones
    \item \textbf{Medical expulsion therapy:} medicines to help stones pass
    \item \textbf{Shock wave lithotripsy:} outpatient treatment, breaking stones with shock waves
    \item \textbf{Ureteroscopy:} endoscopic removal or laser fragmentation of stones
    \item \textbf{Percutaneous nephrolithotomy:} for large or complex stones
    \item \textbf{Stents and drainage:} for obstruction and infection
    \item \textbf{Emergency treatment:} antibiotics and drainage for infected obstruction
    \item \textbf{Preventive therapy:} fluids, dietary advice, and medicines such as thiazides or allopurinol for recurrent stones
\end{itemize}

\section*{Complications}
\begin{itemize}
    \item Obstruction and kidney damage
    \item Sepsis from infected obstruction
    \item Bleeding, infection, and injury from procedures
    \item Retained fragments and recurrence
    \item Stent discomfort and migration
    \item Complications of lithotripsy, including bruising and pain
    \item Repeat treatments when stones are not fully cleared
\end{itemize}

\section*{Prognosis and Outcomes}
Treatment of kidney stones is highly successful, with clearance rates above 90 per cent for most techniques, and modern procedures have made open surgery rare. The main challenge is recurrence, which preventive measures halve. With good hydration, dietary change, and, where needed, medicines, most people avoid further stones, and those that do recur are treated effectively.

\section*{Prevention and Self-Care}
\begin{itemize}
    \item Drink enough fluid to keep urine pale, around 2 to 3 litres daily
    \item Limit salt and reduce animal protein
    \item Follow dietary advice based on your stone type
    \item Have stones analysed and attend metabolic testing
    \item Take preventive medicines as prescribed
    \item Attend follow-up imaging
    \item Seek early care for pain, fever, or difficulty passing urine
    \item Maintain a healthy weight
    \item Stay active
\end{itemize}

\section*{Sources and Further Reading}
The following reputable sources provide further information for healthcare students and patients seeking reliable, current advice about kidney stone treatment.
\begin{itemize}
    \item Urological Society of Australia and New Zealand. \textit{Kidney stone treatment.} https://www.usanz.org.au/
    \item Kidney Health Australia. \textit{Kidney stones and treatment.} https://kidney.org.au/
    \item Healthdirect Australia. \textit{Kidney stones treatment.} https://www.healthdirect.gov.au/
    \item Better Health Channel. \textit{Kidney stones.} https://www.betterhealth.vic.gov.au/
    \item NHS. \textit{Kidney stone treatment.} https://www.nhs.uk/
    \item Mayo Clinic. \textit{Kidney stone treatment.} https://www.mayoclinic.org/
\end{itemize}
